\chapter*{Acknowledgements}
\addcontentsline{toc}{chapter}{Acknowledgements}

This book exists because one person chose to trust another with a dangerous story. The man I call Healer in these pages spent three years teaching me to see what was in front of me. I cannot thank him by name. I hope this book is thanks enough.

Genevieve Prentice designed the ethical framework through which the story finally made sense. Her Dignity Net architecture provided what twenty years of my own thinking could not: a way to tell this story through a moral lens rather than a technical one. Her name is on the cover because she earned it. This book is immeasurably better for her involvement.

My daughters, Gillian and Fiona, grew up in the shadow of this story. They met David when they were six and three. They have lived with a father who could not fully explain what he was working on, or why it mattered, for most of their lives. Their patience and trust mean more than I can say here.

Karen Stephenson raised our daughters through the years this story consumed their father.

Michael Angerman was at the Santa Fe Institute from 1986 through 1995, a DARPA-funded researcher working with Kauffman on neural networks. He went on to Sun Microsystems, where Bill Joy was Chief Scientist. In a conversation in 2019, he confirmed details that no one outside that world could have known, freely and without an NDA to constrain him. He may yet have more to say.

Dave Bannon, physics instructor at Oregon State University, arranged interlibrary loan access to Ivancevic's \textit{Quantum Neural Computation} --- a 929-page Springer monograph that very few libraries hold and fewer people have read. Two six-week loans. That book filled gaps that years of papers could not.

Richard Crandall taught senior quantum mechanics at Reed College with the rigor and joy that made everything after it possible. He died in 2012. I wish he could read this.

Thomas Bury at UC Riverside engaged seriously with the mathematical connections between early warning signals and criticality theory, and his correspondence helped sharpen the empirical foundations of this work.

Lawrence Mysak at McGill provided, unknowingly, a key piece of the mathematical lineage connecting thermohaline circulation models to the same bifurcation structures that appear throughout this book.

Bruce Schneier, Cory Doctorow, and John Gilmore hold copies of this manuscript under a deadman's switch arrangement. Their willingness to serve as custodians of a strange document from a stranger is an act of civic generosity.

% TODO: add Micah Lee, Laura Poitras if contact established before publication

Friends and neighbors in Corvallis lived through these years with me. They know who they are.

Majikan Geoff Wathen trained me in pukulan for ten years. What he taught has nothing to do with quantum computing and everything to do with the kind of person who could write this book.

Jeff Bradford, my college roommate, humbled my ego and taught me how to lose at chess. Both skills proved essential.

Teri Turner taught me about the power of three.

Tanja taught me the value of a good dog.

Magdalena kept me fed and kept the house standing at Pipe House while I wrote the last pieces of this manuscript. She taught me Spanish. Gracias, Magdalena.

\vspace{1cm}

This book was written with substantial assistance from Claude, an AI system built by Anthropic. The full session transcripts are included in the appendix. I used AI as a research partner, structural editor, fact-checker, and sparring partner --- never as a ghostwriter. Every claim, every judgment, every sentence that matters is mine or Genevieve's. The AI helped me find what I already knew. It did not know things I did not. The irony of using an artificial intelligence to write a book about artificial intelligence is not lost on me.

